% META: {"id": "TSPE-DERCONV-CO-022", "chapitre": "TSPE-DERIVATION-CONVEXITE", "type_objet": "corrige", "exercice_ref": "TSPE-DERCONV-EX-022", "capacites_codes": ["C3"], "sources_inspiration": [], "status": "generated"}
% BEGIN-VERIFY
% from sympy import *
% x, y, t = symbols('x y t')
% # Jensen: f(tx+(1-t)y) <= t*f(x) + (1-t)*f(y) pour f convexe
% # pour f=exp, t=1/2:
% lhs = exp((x+y)/2)
% rhs = exp(x)/2 + exp(y)/2
% # verifier avec valeurs numeriques
% assert lhs.subs({x: 0, y: 2}) == exp(1)
% assert rhs.subs({x: 0, y: 2}) == (1 + exp(2))/2
% assert (rhs - lhs).subs({x: 0, y: 2}) > 0
% END-VERIFY
\begin{corrige}{TSPE-DERCONV-EX-022}

\textbf{Question 1.} Par convexite de $\mathrm{e}^x$ et l'inegalite de Jensen :
\[ \mathrm{e}^{\frac{x+y}{2}} \leq \frac{\mathrm{e}^x + \mathrm{e}^y}{2}. \]

\commentaireMarge{Jensen avec $t=\frac12$ : moyenne geometrique $\leq$ moyenne arithmetique (cas exponentielle).}

\textbf{Question 2.} Pour $x=0$ et $y=2$ :
\[ \mathrm{e}^{1} \approx 2{,}718 \quad \text{et} \quad \frac{1 + \mathrm{e}^2}{2} \approx \frac{1 + 7{,}389}{2} \approx 4{,}195. \]
On a bien $\mathrm{e} < \dfrac{1 + \mathrm{e}^2}{2}$.

\end{corrige}
